% Preambul comun pentru toate handout-urile training-ului.
% Fiecare fisier .tex face % Preambul comun pentru toate handout-urile training-ului.
% Fiecare fisier .tex face % Preambul comun pentru toate handout-urile training-ului.
% Fiecare fisier .tex face % Preambul comun pentru toate handout-urile training-ului.
% Fiecare fisier .tex face \input{preamble} dupa \documentclass.

\usepackage[utf8]{inputenc}
\usepackage[T1]{fontenc}
% Body font: Lato -- modern, lizibil, profesional. Fallback la TeX Gyre Heros
% (Helvetica-like, mereu disponibil) daca Lato nu e instalat in MiKTeX.
\IfFileExists{lato.sty}{%
  \usepackage[default]{lato}%
}{%
  \usepackage{tgheros}%
  \renewcommand{\familydefault}{\sfdefault}%
}
\usepackage[romanian]{babel}
\usepackage[a4paper, margin=2.5cm, headheight=16pt, top=2.4cm, bottom=2.4cm]{geometry}
\usepackage[final, protrusion=true, expansion=false]{microtype}
\usepackage{setspace}
\usepackage{xcolor}
\usepackage{amssymb}
\usepackage{listings}
\usepackage{hyperref}
\usepackage{enumitem}
\usepackage{titlesec}
\usepackage{fancyhdr}
\usepackage{parskip}
\usepackage[most]{tcolorbox}
\usepackage{tocloft}
\usepackage{graphicx}
\usepackage{tikz}
\usepackage{fontawesome5}
\usepackage{lettrine}
\usetikzlibrary{shadows.blur, shapes.geometric, positioning, calc,
  decorations.pathmorphing, patterns, fadings, backgrounds}

% =========================================================
%   PALETA DE CULORI
% =========================================================

% Paleta finala (3 BEST + 3 derivate in acelasi ton)
\definecolor{themePrereq}{HTML}{5C6B7A}    % slate (pre-rechizite, neutral)
\definecolor{themeHtml}{HTML}{FAA51E}      % BEST orange
\definecolor{themeCss}{HTML}{0673BA}       % BEST blue
\definecolor{themeJs}{HTML}{ECC100}        % gold (analog cu orange, distinct)
\definecolor{themeDeploy}{HTML}{75C044}    % BEST green
\definecolor{themeSeo}{HTML}{B344C9}       % violet (analog cu blue)
\definecolor{themeLighthouse}{HTML}{DD4444}% coral red (complementar cu green)
\definecolor{themeResurse}{HTML}{5C6B7A}   % slate (resurse, neutral)
\definecolor{themeBonus}{HTML}{7E22CE}     % deep purple (bonus / avansat, self-study only)
% Aliasuri pentru compatibilitate (folosit in main.tex existent)
\colorlet{themeIntro}{themeHtml}
\colorlet{themeExercitiu}{themeLighthouse}

% Background pagina si elemente -- calduros
\definecolor{pageBg}{HTML}{FFFCF6}         % crema foarte subtila
\definecolor{accentSoft}{HTML}{F4ECDD}     % accent cald
\definecolor{textDim}{HTML}{64748B}        % gri pentru text secundar

% Culoarea curenta a partii (suprascrisa de fiecare fisier prin \setpart)
\colorlet{partcolor}{themeIntro}

% Iconita partii curente (suprascrisa prin \setpart)
\newcommand{\particon}{\faStar}
\newcommand{\partlabel}{Training Web}

% \setpart{culoare}{iconita}
\newcommand{\setpart}[2]{%
  \colorlet{partcolor}{#1}%
  \renewcommand{\particon}{#2}%
}
% \setpart{culoare}{iconita}{eticheta} -- versiunea cu eticheta custom in header
\newcommand{\setpartlabel}[3]{%
  \colorlet{partcolor}{#1}%
  \renewcommand{\particon}{#2}%
  \renewcommand{\partlabel}{#3}%
}

% Culori pentru callout boxes
\definecolor{noteBg}{HTML}{FEF9C3}        \definecolor{noteFg}{HTML}{A16207}
\definecolor{tipBg}{HTML}{DBEAFE}         \definecolor{tipFg}{HTML}{1D4ED8}
\definecolor{warnBg}{HTML}{FEE2E2}        \definecolor{warnFg}{HTML}{B91C1C}
\definecolor{gameBg}{HTML}{FFEDD5}        \definecolor{gameFg}{HTML}{C2410C}
\definecolor{recapBg}{HTML}{DCFCE7}       \definecolor{recapFg}{HTML}{15803D}
\definecolor{exerBg}{HTML}{F3E8FF}        \definecolor{exerFg}{HTML}{7C3AED}
\definecolor{triviaBg}{HTML}{FCE7F3}      \definecolor{triviaFg}{HTML}{BE185D}

% Culori pentru cod -- paleta "GitHub Light" (calm, profesional)
\definecolor{codeBg}{HTML}{F6F8FA}
\definecolor{codeBorder}{HTML}{D0D7DE}
\definecolor{codeComment}{HTML}{6E7781}
\definecolor{codeKeyword}{HTML}{CF222E}
\definecolor{codeString}{HTML}{0A3069}
\definecolor{codeTag}{HTML}{116329}

% =========================================================
%   PAGINA & LIZIBILITATE (warmth)
% =========================================================
\pagecolor{pageBg}
\setstretch{1.18}
\setlength{\parskip}{0.5em}

% =========================================================
%   CUPRINS (TOC)
% =========================================================
\setcounter{tocdepth}{0}         % numai entry-uri \part (capitole), fara sectiuni
\renewcommand{\contentsname}{Cuprins}

% Stil cuprins -- numai capitolele, asezate elegant
\renewcommand{\cftpartfont}{\sffamily\bfseries\color{themeText}}
\renewcommand{\cftpartpagefont}{\sffamily\bfseries\color{themeText}}
\renewcommand{\cftpartleader}{\hfill}
\renewcommand{\cftpartafterpnum}{\par\addvspace{2pt}}
\setlength{\cftbeforepartskip}{4pt}

% =========================================================
%   HYPERREF
% =========================================================
\hypersetup{
  colorlinks=true,
  linkcolor=partcolor,
  urlcolor=partcolor,
  citecolor=partcolor,
  pdfborderstyle={/S/U/W 1}
}

% =========================================================
%   HEADER & FOOTER -- minimaliste
% =========================================================
\pagestyle{fancy}
\fancyhf{}
\fancyhead[L]{\footnotesize\sffamily\color{black!55}\partlabel}
\fancyhead[R]{\footnotesize\sffamily\color{black!55}BEST Ia\textcommabelow{s}i}
\fancyfoot[C]{\small\sffamily\color{partcolor}\textbf{\thepage}}
\renewcommand{\headrulewidth}{0pt}
\renewcommand{\footrulewidth}{0pt}

% =========================================================
%   STIL TITLURI -- sobru, stiintific
% =========================================================
\titleformat{\section}
  {\Large\sffamily\bfseries\color{themeText}}
  {\color{partcolor}\sffamily\bfseries\thesection.\hspace{0.5em}}
  {0pt}{}
\titlespacing*{\section}{0pt}{1.4em}{0.4em}

\titleformat{\subsection}
  {\normalsize\sffamily\bfseries\color{black!75}}
  {\color{partcolor}\thesubsection\hspace{0.5em}}{0pt}{}
\titlespacing*{\subsection}{0pt}{0.8em}{0.2em}

\titleformat{\subsubsection}
  {\normalsize\sffamily\bfseries\color{black!70}}
  {}{0pt}{}

\definecolor{themeText}{HTML}{1A1A1A}

% =========================================================
%   LISTINGS -- COD
% =========================================================
\lstdefinestyle{base}{
  backgroundcolor=\color{codeBg},
  basicstyle=\ttfamily\small,
  commentstyle=\color{codeComment}\itshape,
  keywordstyle=\color{codeKeyword}\bfseries,
  stringstyle=\color{codeString},
  numberstyle=\tiny\color{codeComment},
  numbers=left,
  numbersep=8pt,
  frame=leftline,
  framerule=2.5pt,
  rulecolor=\color{partcolor},
  framexleftmargin=4pt,
  breaklines=true,
  breakatwhitespace=true,
  showstringspaces=false,
  tabsize=2,
  xleftmargin=14pt,
  xrightmargin=4pt,
  aboveskip=10pt,
  belowskip=10pt,
  literate=
    {ă}{{\u a}}1 {Ă}{{\u A}}1
    {â}{{\^a}}1 {Â}{{\^A}}1
    {î}{{\^\i}}1 {Î}{{\^I}}1
    {ș}{{\textcommabelow s}}1 {Ș}{{\textcommabelow S}}1
    {ț}{{\textcommabelow t}}1 {Ț}{{\textcommabelow T}}1
}

\lstdefinelanguage{HTML5}{
  morekeywords={html,head,body,title,meta,link,script,style,
    h1,h2,h3,h4,h5,h6,p,a,img,div,span,
    header,main,section,article,aside,nav,footer,
    ul,ol,li,button,input,form,label,
    strong,em,br,hr,DOCTYPE},
  sensitive=true,
  morecomment=[s]{<!--}{-->},
  morestring=[b]",
  morestring=[b]'
}
\lstdefinestyle{html}{style=base, language=HTML5,
  keywordstyle=\color{codeTag}\bfseries}

\lstdefinelanguage{CSSLang}{
  morekeywords={display,flex,grid,block,inline,none,
    color,background,background-color,
    padding,padding-top,padding-right,padding-bottom,padding-left,
    margin,margin-top,margin-right,margin-bottom,margin-left,
    border,border-radius,border-color,outline,
    width,max-width,min-width,height,max-height,min-height,
    font-family,font-size,font-weight,line-height,text-align,
    text-decoration,letter-spacing,
    flex-direction,justify-content,align-items,gap,flex-wrap,flex-grow,
    grid-template-columns,grid-template-rows,grid-column,grid-row,
    position,top,right,bottom,left,z-index,
    box-sizing,opacity,transition,transform,
    cursor,overflow,visibility,
    media,keyframes,from,to,important,root,hover,focus,active,
    auto,solid,dashed,dotted,bold,normal,center,stretch,
    space-between,space-around,flex-start,flex-end,column,row,wrap},
  sensitive=true,
  morecomment=[s]{/*}{*/},
  morestring=[b]",
  morestring=[b]'
}
\lstdefinestyle{css}{style=base, language=CSSLang,
  keywordstyle=\color{codeKeyword}\bfseries}

\lstdefinelanguage{JSLang}{
  morekeywords={let,const,var,function,return,if,else,for,while,
    true,false,null,undefined,new,this,class,extends,
    document,console,window,
    addEventListener,querySelector,querySelectorAll,
    getElementById,getElementsByClassName,
    classList,toggle,add,remove,contains,
    style,textContent,innerHTML,value,
    log,warn,error},
  sensitive=true,
  morecomment=[l]{//},
  morecomment=[s]{/*}{*/},
  morestring=[b]",
  morestring=[b]',
  morestring=[b]`
}
\lstdefinestyle{js}{style=base, language=JSLang,
  keywordstyle=\color{codeKeyword}\bfseries}

\lstdefinelanguage{BashLang}{
  morekeywords={git,cd,mkdir,ls,code,echo,cat,curl,wget,sudo,
    init,add,commit,push,pull,clone,remote,branch,checkout,merge,status,log,
    config,fetch,rebase,reset,stash,diff,tag,
    install,update,run,build,start,test},
  sensitive=true,
  morecomment=[l]{\#},
  morestring=[b]",
  morestring=[b]'
}
\lstdefinestyle{bash}{style=base, language=BashLang,
  keywordstyle=\color{codeKeyword}\bfseries}

% Inline code -- mai discret
\newcommand{\code}[1]{\textcolor{partcolor!85!black}{\ttfamily\small #1}}

% =========================================================
%   CALLOUT BOXES -- minimaliste (bara stanga subtila + label discret)
% =========================================================
\tcbset{
  calloutbase/.style={
    enhanced,
    breakable,
    sharp corners,
    boxrule=0pt,
    leftrule=2.5pt,
    arc=0pt,
    colback=pageBg,
    fonttitle=\sffamily\small\bfseries,
    coltitle=tcbcolframe,
    attach title to upper={\par\vspace{2pt}},
    title after break={},
    before skip=10pt,
    after skip=10pt,
    left=12pt, right=4pt, top=4pt, bottom=4pt,
    fontupper=\small,
  }
}

\newtcolorbox{nota}[1][NOT\u A]{
  calloutbase,
  colframe=noteFg,
  title={\faInfoCircle\ \uppercase\expandafter{#1}},
}
\newtcolorbox{ponturi}[1][PONT]{
  calloutbase,
  colframe=tipFg,
  title={\faLightbulb[regular]\ \uppercase\expandafter{#1}},
}
\newtcolorbox{atentie}[1][ATEN\textcommabelow{T}IE]{
  calloutbase,
  colframe=warnFg,
  title={\faExclamationTriangle\ \uppercase\expandafter{#1}},
}
\newtcolorbox{joc}[1][APLICA\textcommabelow{T}IE PRACTIC\u A]{
  calloutbase,
  colframe=partcolor,
  title={\faGamepad\ \uppercase\expandafter{#1}},
}
\newtcolorbox{recap}[1][RECAP]{
  calloutbase,
  colframe=recapFg,
  title={\faCheckCircle\ \uppercase\expandafter{#1}},
}
\newtcolorbox{curiozitate}[1][\textcommabelow{S}TIAI C\u A?]{
  calloutbase,
  colframe=triviaFg,
  title={\faMagic\ \uppercase\expandafter{#1}},
}
\newtcolorbox{exercitiu}[1][EXERCI\textcommabelow{T}IU]{
  calloutbase,
  colframe=exerFg,
  title={\faPen\ \uppercase\expandafter{#1}},
}

% --- Outcome / Why / Self-check (structura tutorial) ----
\definecolor{outBg}{HTML}{E0F2FE}        \definecolor{outFg}{HTML}{0369A1}
\definecolor{whyBg}{HTML}{FEF3C7}        \definecolor{whyFg}{HTML}{B45309}
\definecolor{checkBg}{HTML}{F3E8FF}      \definecolor{checkFg}{HTML}{7C3AED}

\newtcolorbox{outcomes}[1][CE INVE\textcommabelow{T}I]{
  calloutbase,
  colframe=outFg,
  title={\faBullseye\ \uppercase\expandafter{#1}},
}
\newtcolorbox{whymatters}[1][DE CE]{
  calloutbase,
  colframe=whyFg,
  title={\faQuestionCircle\ \uppercase\expandafter{#1}},
}
\newtcolorbox{selfcheck}[1][PROVOCARE]{
  calloutbase,
  colframe=checkFg,
  title={\faTasks\ \uppercase\expandafter{#1}},
}

% --- Bonus / Avansat (continut auto-studiu, distinct vizual) ----
% Diferit fata de celelalte callout-uri: fundal usor colorat (nu pageBg),
% bara stanga mai groasa, icon racheta. Beginners vad clar si pot sari peste.
\definecolor{bonusBg}{HTML}{F6EEFB}        \definecolor{bonusFg}{HTML}{6B21A8}

\newtcolorbox{bonusbox}[1][BONUS \,\textbullet\, AVANSAT]{
  enhanced,
  breakable,
  sharp corners,
  boxrule=0pt,
  leftrule=4pt,
  arc=0pt,
  colback=bonusBg,
  colframe=themeBonus,
  fonttitle=\sffamily\small\bfseries,
  coltitle=bonusFg,
  attach title to upper={\par\vspace{2pt}},
  title after break={\faRocket\ \textbf{\color{bonusFg}continuare bonus}},
  before skip=14pt,
  after skip=14pt,
  left=14pt, right=6pt, top=6pt, bottom=6pt,
  fontupper=\small,
  title={\faRocket\ \uppercase\expandafter{#1}},
}

% Bridge la finalul fiecarui capitol
\newcommand{\bridge}[1]{%
  \par\vspace{4pt}%
  \begin{flushleft}%
    {\sffamily\small\color{partcolor}$\rightarrow$\ \textbf{Urmeaza:}}\quad{\small\color{black!70}#1}%
  \end{flushleft}%
}

% --- Asa da / Asa nu (compare bune vs rele practici) ----
\definecolor{daBg}{HTML}{DCFCE7}        \definecolor{daFg}{HTML}{15803D}
\definecolor{nuBg}{HTML}{FEE2E2}        \definecolor{nuFg}{HTML}{B91C1C}

\newtcolorbox{dado}[1][A\textcommabelow{S}A DA]{
  calloutbase,
  colframe=daFg,
  title={\faCheckCircle\ \uppercase\expandafter{#1}},
}
\newtcolorbox{dano}[1][A\textcommabelow{S}A NU]{
  calloutbase,
  colframe=nuFg,
  title={\faTimesCircle\ \uppercase\expandafter{#1}},
}

% =========================================================
%   BANNER MINIMALIST (\handout{titlu}{subtitlu})
%   Eticheta sus, titlu mare, subtitlu, linie subtila colorata.
% =========================================================
\newcommand{\handout}[2]{%
  \noindent\begin{minipage}{0.04\linewidth}\color{partcolor}\rule{4pt}{2.6em}\end{minipage}%
  \begin{minipage}{0.94\linewidth}
    {\sffamily\fontsize{30}{34}\selectfont\textbf{\color{themeText}#1}}\\[4pt]
    {\sffamily\large\color{black!55}#2}
  \end{minipage}
  \par\vspace{14pt}
  {\color{partcolor!40}\hrule height 0.5pt}
  \par\vspace{14pt}
}

% =========================================================
%   BULLETS COLORATI (itemize)
% =========================================================
\setlist[itemize,1]{label={\color{partcolor}\textbf{\textbullet}}}
\setlist[itemize,2]{label={\color{partcolor!70}\textbf{\textendash}}}
\setlist[enumerate,1]{label={\color{partcolor}\textbf{\arabic*.}}}

% =========================================================
%   PULL QUOTE / ACCENT
% =========================================================
\newcommand{\pullquote}[1]{%
  \begin{tcolorbox}[
    enhanced,
    colback=partcolor!8,
    colframe=partcolor,
    boxrule=0pt,
    leftrule=4pt,
    sharp corners,
    arc=0pt,
    fontupper=\itshape\large,
    left=14pt, right=14pt, top=8pt, bottom=8pt,
    overlay={%
      \node[partcolor!40, font=\fontsize{40}{40}\selectfont, anchor=north west]
        at ([xshift=2pt,yshift=2pt]frame.north west) {\textbf{``}};
    },
  ]
    #1
  \end{tcolorbox}
}

% =========================================================
%   DIVIDER DECORATIV (foloseste \partdivider in document)
%   Atentie: NU folosi numele "\sectionbreak" -- e rezervat
%   intern de titlesec si va aparea automat intre sectiuni.
% =========================================================
\newcommand{\partdivider}{%
  \par\vspace{0.6em}
  \begin{center}
    \tikz[baseline=-0.5ex]{
      \fill[partcolor] (0,0) circle (2pt);
      \fill[partcolor!60] (10pt,0) circle (2pt);
      \fill[partcolor!30] (20pt,0) circle (2pt);
    }
  \end{center}
  \vspace{0.3em}
}

% =========================================================
%   BOX MODEL DIAGRAM (TikZ)
% =========================================================
\newcommand{\boxmodelfigure}{%
\begin{center}
\begin{tikzpicture}[scale=1, transform shape, every node/.style={font=\sffamily\small}]
  % Margin (outermost)
  \fill[orange!25, rounded corners=4pt] (0,0) rectangle (8,5);
  \node[anchor=north west] at (0.15, 4.85) {\textbf{\color{orange!70!black}\faExpand\ margin}};

  % Border
  \fill[gray!40, rounded corners=3pt] (1,0.7) rectangle (7, 4.3);
  \draw[ultra thick, gray!70!black, rounded corners=3pt] (1,0.7) rectangle (7,4.3);
  \node[anchor=north west] at (1.1, 4.2) {\textbf{\color{gray!50!black}\faSquare\ border}};

  % Padding
  \fill[green!25, rounded corners=2pt] (1.3, 1) rectangle (6.7, 4);
  \node[anchor=north west] at (1.4, 3.9) {\textbf{\color{green!50!black}\faCompress\ padding}};

  % Content (innermost)
  \fill[blue!25, rounded corners=2pt] (2.2, 1.7) rectangle (5.8, 3.3);
  \node[anchor=north west] at (2.3, 3.2) {\textbf{\color{blue!60!black}\faFile*\ content}};
  \node at (4, 2.5) {\textit{textul, imaginea, ce e inauntru}};
\end{tikzpicture}
\end{center}
}

% =========================================================
%   CARD "ICON + TEXT" (utilitate pentru handout-uri)
% =========================================================
\newcommand{\iconcard}[3]{%
  \begin{tcolorbox}[
    enhanced,
    colback=partcolor!8,
    colframe=partcolor!30,
    boxrule=0.5pt,
    arc=6pt,
    left=12pt, right=12pt, top=8pt, bottom=8pt,
  ]
  \begin{minipage}[t]{0.08\linewidth}
    \color{partcolor}\fontsize{20}{20}\selectfont #1
  \end{minipage}%
  \begin{minipage}[t]{0.9\linewidth}
    \textbf{\color{partcolor!85!black}#2}\\
    #3
  \end{minipage}
  \end{tcolorbox}%
}
 dupa \documentclass.

\usepackage[utf8]{inputenc}
\usepackage[T1]{fontenc}
% Body font: Lato -- modern, lizibil, profesional. Fallback la TeX Gyre Heros
% (Helvetica-like, mereu disponibil) daca Lato nu e instalat in MiKTeX.
\IfFileExists{lato.sty}{%
  \usepackage[default]{lato}%
}{%
  \usepackage{tgheros}%
  \renewcommand{\familydefault}{\sfdefault}%
}
\usepackage[romanian]{babel}
\usepackage[a4paper, margin=2.5cm, headheight=16pt, top=2.4cm, bottom=2.4cm]{geometry}
\usepackage[final, protrusion=true, expansion=false]{microtype}
\usepackage{setspace}
\usepackage{xcolor}
\usepackage{amssymb}
\usepackage{listings}
\usepackage{hyperref}
\usepackage{enumitem}
\usepackage{titlesec}
\usepackage{fancyhdr}
\usepackage{parskip}
\usepackage[most]{tcolorbox}
\usepackage{tocloft}
\usepackage{graphicx}
\usepackage{tikz}
\usepackage{fontawesome5}
\usepackage{lettrine}
\usetikzlibrary{shadows.blur, shapes.geometric, positioning, calc,
  decorations.pathmorphing, patterns, fadings, backgrounds}

% =========================================================
%   PALETA DE CULORI
% =========================================================

% Paleta finala (3 BEST + 3 derivate in acelasi ton)
\definecolor{themePrereq}{HTML}{5C6B7A}    % slate (pre-rechizite, neutral)
\definecolor{themeHtml}{HTML}{FAA51E}      % BEST orange
\definecolor{themeCss}{HTML}{0673BA}       % BEST blue
\definecolor{themeJs}{HTML}{ECC100}        % gold (analog cu orange, distinct)
\definecolor{themeDeploy}{HTML}{75C044}    % BEST green
\definecolor{themeSeo}{HTML}{B344C9}       % violet (analog cu blue)
\definecolor{themeLighthouse}{HTML}{DD4444}% coral red (complementar cu green)
\definecolor{themeResurse}{HTML}{5C6B7A}   % slate (resurse, neutral)
\definecolor{themeBonus}{HTML}{7E22CE}     % deep purple (bonus / avansat, self-study only)
% Aliasuri pentru compatibilitate (folosit in main.tex existent)
\colorlet{themeIntro}{themeHtml}
\colorlet{themeExercitiu}{themeLighthouse}

% Background pagina si elemente -- calduros
\definecolor{pageBg}{HTML}{FFFCF6}         % crema foarte subtila
\definecolor{accentSoft}{HTML}{F4ECDD}     % accent cald
\definecolor{textDim}{HTML}{64748B}        % gri pentru text secundar

% Culoarea curenta a partii (suprascrisa de fiecare fisier prin \setpart)
\colorlet{partcolor}{themeIntro}

% Iconita partii curente (suprascrisa prin \setpart)
\newcommand{\particon}{\faStar}
\newcommand{\partlabel}{Training Web}

% \setpart{culoare}{iconita}
\newcommand{\setpart}[2]{%
  \colorlet{partcolor}{#1}%
  \renewcommand{\particon}{#2}%
}
% \setpart{culoare}{iconita}{eticheta} -- versiunea cu eticheta custom in header
\newcommand{\setpartlabel}[3]{%
  \colorlet{partcolor}{#1}%
  \renewcommand{\particon}{#2}%
  \renewcommand{\partlabel}{#3}%
}

% Culori pentru callout boxes
\definecolor{noteBg}{HTML}{FEF9C3}        \definecolor{noteFg}{HTML}{A16207}
\definecolor{tipBg}{HTML}{DBEAFE}         \definecolor{tipFg}{HTML}{1D4ED8}
\definecolor{warnBg}{HTML}{FEE2E2}        \definecolor{warnFg}{HTML}{B91C1C}
\definecolor{gameBg}{HTML}{FFEDD5}        \definecolor{gameFg}{HTML}{C2410C}
\definecolor{recapBg}{HTML}{DCFCE7}       \definecolor{recapFg}{HTML}{15803D}
\definecolor{exerBg}{HTML}{F3E8FF}        \definecolor{exerFg}{HTML}{7C3AED}
\definecolor{triviaBg}{HTML}{FCE7F3}      \definecolor{triviaFg}{HTML}{BE185D}

% Culori pentru cod -- paleta "GitHub Light" (calm, profesional)
\definecolor{codeBg}{HTML}{F6F8FA}
\definecolor{codeBorder}{HTML}{D0D7DE}
\definecolor{codeComment}{HTML}{6E7781}
\definecolor{codeKeyword}{HTML}{CF222E}
\definecolor{codeString}{HTML}{0A3069}
\definecolor{codeTag}{HTML}{116329}

% =========================================================
%   PAGINA & LIZIBILITATE (warmth)
% =========================================================
\pagecolor{pageBg}
\setstretch{1.18}
\setlength{\parskip}{0.5em}

% =========================================================
%   CUPRINS (TOC)
% =========================================================
\setcounter{tocdepth}{0}         % numai entry-uri \part (capitole), fara sectiuni
\renewcommand{\contentsname}{Cuprins}

% Stil cuprins -- numai capitolele, asezate elegant
\renewcommand{\cftpartfont}{\sffamily\bfseries\color{themeText}}
\renewcommand{\cftpartpagefont}{\sffamily\bfseries\color{themeText}}
\renewcommand{\cftpartleader}{\hfill}
\renewcommand{\cftpartafterpnum}{\par\addvspace{2pt}}
\setlength{\cftbeforepartskip}{4pt}

% =========================================================
%   HYPERREF
% =========================================================
\hypersetup{
  colorlinks=true,
  linkcolor=partcolor,
  urlcolor=partcolor,
  citecolor=partcolor,
  pdfborderstyle={/S/U/W 1}
}

% =========================================================
%   HEADER & FOOTER -- minimaliste
% =========================================================
\pagestyle{fancy}
\fancyhf{}
\fancyhead[L]{\footnotesize\sffamily\color{black!55}\partlabel}
\fancyhead[R]{\footnotesize\sffamily\color{black!55}BEST Ia\textcommabelow{s}i}
\fancyfoot[C]{\small\sffamily\color{partcolor}\textbf{\thepage}}
\renewcommand{\headrulewidth}{0pt}
\renewcommand{\footrulewidth}{0pt}

% =========================================================
%   STIL TITLURI -- sobru, stiintific
% =========================================================
\titleformat{\section}
  {\Large\sffamily\bfseries\color{themeText}}
  {\color{partcolor}\sffamily\bfseries\thesection.\hspace{0.5em}}
  {0pt}{}
\titlespacing*{\section}{0pt}{1.4em}{0.4em}

\titleformat{\subsection}
  {\normalsize\sffamily\bfseries\color{black!75}}
  {\color{partcolor}\thesubsection\hspace{0.5em}}{0pt}{}
\titlespacing*{\subsection}{0pt}{0.8em}{0.2em}

\titleformat{\subsubsection}
  {\normalsize\sffamily\bfseries\color{black!70}}
  {}{0pt}{}

\definecolor{themeText}{HTML}{1A1A1A}

% =========================================================
%   LISTINGS -- COD
% =========================================================
\lstdefinestyle{base}{
  backgroundcolor=\color{codeBg},
  basicstyle=\ttfamily\small,
  commentstyle=\color{codeComment}\itshape,
  keywordstyle=\color{codeKeyword}\bfseries,
  stringstyle=\color{codeString},
  numberstyle=\tiny\color{codeComment},
  numbers=left,
  numbersep=8pt,
  frame=leftline,
  framerule=2.5pt,
  rulecolor=\color{partcolor},
  framexleftmargin=4pt,
  breaklines=true,
  breakatwhitespace=true,
  showstringspaces=false,
  tabsize=2,
  xleftmargin=14pt,
  xrightmargin=4pt,
  aboveskip=10pt,
  belowskip=10pt,
  literate=
    {ă}{{\u a}}1 {Ă}{{\u A}}1
    {â}{{\^a}}1 {Â}{{\^A}}1
    {î}{{\^\i}}1 {Î}{{\^I}}1
    {ș}{{\textcommabelow s}}1 {Ș}{{\textcommabelow S}}1
    {ț}{{\textcommabelow t}}1 {Ț}{{\textcommabelow T}}1
}

\lstdefinelanguage{HTML5}{
  morekeywords={html,head,body,title,meta,link,script,style,
    h1,h2,h3,h4,h5,h6,p,a,img,div,span,
    header,main,section,article,aside,nav,footer,
    ul,ol,li,button,input,form,label,
    strong,em,br,hr,DOCTYPE},
  sensitive=true,
  morecomment=[s]{<!--}{-->},
  morestring=[b]",
  morestring=[b]'
}
\lstdefinestyle{html}{style=base, language=HTML5,
  keywordstyle=\color{codeTag}\bfseries}

\lstdefinelanguage{CSSLang}{
  morekeywords={display,flex,grid,block,inline,none,
    color,background,background-color,
    padding,padding-top,padding-right,padding-bottom,padding-left,
    margin,margin-top,margin-right,margin-bottom,margin-left,
    border,border-radius,border-color,outline,
    width,max-width,min-width,height,max-height,min-height,
    font-family,font-size,font-weight,line-height,text-align,
    text-decoration,letter-spacing,
    flex-direction,justify-content,align-items,gap,flex-wrap,flex-grow,
    grid-template-columns,grid-template-rows,grid-column,grid-row,
    position,top,right,bottom,left,z-index,
    box-sizing,opacity,transition,transform,
    cursor,overflow,visibility,
    media,keyframes,from,to,important,root,hover,focus,active,
    auto,solid,dashed,dotted,bold,normal,center,stretch,
    space-between,space-around,flex-start,flex-end,column,row,wrap},
  sensitive=true,
  morecomment=[s]{/*}{*/},
  morestring=[b]",
  morestring=[b]'
}
\lstdefinestyle{css}{style=base, language=CSSLang,
  keywordstyle=\color{codeKeyword}\bfseries}

\lstdefinelanguage{JSLang}{
  morekeywords={let,const,var,function,return,if,else,for,while,
    true,false,null,undefined,new,this,class,extends,
    document,console,window,
    addEventListener,querySelector,querySelectorAll,
    getElementById,getElementsByClassName,
    classList,toggle,add,remove,contains,
    style,textContent,innerHTML,value,
    log,warn,error},
  sensitive=true,
  morecomment=[l]{//},
  morecomment=[s]{/*}{*/},
  morestring=[b]",
  morestring=[b]',
  morestring=[b]`
}
\lstdefinestyle{js}{style=base, language=JSLang,
  keywordstyle=\color{codeKeyword}\bfseries}

\lstdefinelanguage{BashLang}{
  morekeywords={git,cd,mkdir,ls,code,echo,cat,curl,wget,sudo,
    init,add,commit,push,pull,clone,remote,branch,checkout,merge,status,log,
    config,fetch,rebase,reset,stash,diff,tag,
    install,update,run,build,start,test},
  sensitive=true,
  morecomment=[l]{\#},
  morestring=[b]",
  morestring=[b]'
}
\lstdefinestyle{bash}{style=base, language=BashLang,
  keywordstyle=\color{codeKeyword}\bfseries}

% Inline code -- mai discret
\newcommand{\code}[1]{\textcolor{partcolor!85!black}{\ttfamily\small #1}}

% =========================================================
%   CALLOUT BOXES -- minimaliste (bara stanga subtila + label discret)
% =========================================================
\tcbset{
  calloutbase/.style={
    enhanced,
    breakable,
    sharp corners,
    boxrule=0pt,
    leftrule=2.5pt,
    arc=0pt,
    colback=pageBg,
    fonttitle=\sffamily\small\bfseries,
    coltitle=tcbcolframe,
    attach title to upper={\par\vspace{2pt}},
    title after break={},
    before skip=10pt,
    after skip=10pt,
    left=12pt, right=4pt, top=4pt, bottom=4pt,
    fontupper=\small,
  }
}

\newtcolorbox{nota}[1][NOT\u A]{
  calloutbase,
  colframe=noteFg,
  title={\faInfoCircle\ \uppercase\expandafter{#1}},
}
\newtcolorbox{ponturi}[1][PONT]{
  calloutbase,
  colframe=tipFg,
  title={\faLightbulb[regular]\ \uppercase\expandafter{#1}},
}
\newtcolorbox{atentie}[1][ATEN\textcommabelow{T}IE]{
  calloutbase,
  colframe=warnFg,
  title={\faExclamationTriangle\ \uppercase\expandafter{#1}},
}
\newtcolorbox{joc}[1][APLICA\textcommabelow{T}IE PRACTIC\u A]{
  calloutbase,
  colframe=partcolor,
  title={\faGamepad\ \uppercase\expandafter{#1}},
}
\newtcolorbox{recap}[1][RECAP]{
  calloutbase,
  colframe=recapFg,
  title={\faCheckCircle\ \uppercase\expandafter{#1}},
}
\newtcolorbox{curiozitate}[1][\textcommabelow{S}TIAI C\u A?]{
  calloutbase,
  colframe=triviaFg,
  title={\faMagic\ \uppercase\expandafter{#1}},
}
\newtcolorbox{exercitiu}[1][EXERCI\textcommabelow{T}IU]{
  calloutbase,
  colframe=exerFg,
  title={\faPen\ \uppercase\expandafter{#1}},
}

% --- Outcome / Why / Self-check (structura tutorial) ----
\definecolor{outBg}{HTML}{E0F2FE}        \definecolor{outFg}{HTML}{0369A1}
\definecolor{whyBg}{HTML}{FEF3C7}        \definecolor{whyFg}{HTML}{B45309}
\definecolor{checkBg}{HTML}{F3E8FF}      \definecolor{checkFg}{HTML}{7C3AED}

\newtcolorbox{outcomes}[1][CE INVE\textcommabelow{T}I]{
  calloutbase,
  colframe=outFg,
  title={\faBullseye\ \uppercase\expandafter{#1}},
}
\newtcolorbox{whymatters}[1][DE CE]{
  calloutbase,
  colframe=whyFg,
  title={\faQuestionCircle\ \uppercase\expandafter{#1}},
}
\newtcolorbox{selfcheck}[1][PROVOCARE]{
  calloutbase,
  colframe=checkFg,
  title={\faTasks\ \uppercase\expandafter{#1}},
}

% --- Bonus / Avansat (continut auto-studiu, distinct vizual) ----
% Diferit fata de celelalte callout-uri: fundal usor colorat (nu pageBg),
% bara stanga mai groasa, icon racheta. Beginners vad clar si pot sari peste.
\definecolor{bonusBg}{HTML}{F6EEFB}        \definecolor{bonusFg}{HTML}{6B21A8}

\newtcolorbox{bonusbox}[1][BONUS \,\textbullet\, AVANSAT]{
  enhanced,
  breakable,
  sharp corners,
  boxrule=0pt,
  leftrule=4pt,
  arc=0pt,
  colback=bonusBg,
  colframe=themeBonus,
  fonttitle=\sffamily\small\bfseries,
  coltitle=bonusFg,
  attach title to upper={\par\vspace{2pt}},
  title after break={\faRocket\ \textbf{\color{bonusFg}continuare bonus}},
  before skip=14pt,
  after skip=14pt,
  left=14pt, right=6pt, top=6pt, bottom=6pt,
  fontupper=\small,
  title={\faRocket\ \uppercase\expandafter{#1}},
}

% Bridge la finalul fiecarui capitol
\newcommand{\bridge}[1]{%
  \par\vspace{4pt}%
  \begin{flushleft}%
    {\sffamily\small\color{partcolor}$\rightarrow$\ \textbf{Urmeaza:}}\quad{\small\color{black!70}#1}%
  \end{flushleft}%
}

% --- Asa da / Asa nu (compare bune vs rele practici) ----
\definecolor{daBg}{HTML}{DCFCE7}        \definecolor{daFg}{HTML}{15803D}
\definecolor{nuBg}{HTML}{FEE2E2}        \definecolor{nuFg}{HTML}{B91C1C}

\newtcolorbox{dado}[1][A\textcommabelow{S}A DA]{
  calloutbase,
  colframe=daFg,
  title={\faCheckCircle\ \uppercase\expandafter{#1}},
}
\newtcolorbox{dano}[1][A\textcommabelow{S}A NU]{
  calloutbase,
  colframe=nuFg,
  title={\faTimesCircle\ \uppercase\expandafter{#1}},
}

% =========================================================
%   BANNER MINIMALIST (\handout{titlu}{subtitlu})
%   Eticheta sus, titlu mare, subtitlu, linie subtila colorata.
% =========================================================
\newcommand{\handout}[2]{%
  \noindent\begin{minipage}{0.04\linewidth}\color{partcolor}\rule{4pt}{2.6em}\end{minipage}%
  \begin{minipage}{0.94\linewidth}
    {\sffamily\fontsize{30}{34}\selectfont\textbf{\color{themeText}#1}}\\[4pt]
    {\sffamily\large\color{black!55}#2}
  \end{minipage}
  \par\vspace{14pt}
  {\color{partcolor!40}\hrule height 0.5pt}
  \par\vspace{14pt}
}

% =========================================================
%   BULLETS COLORATI (itemize)
% =========================================================
\setlist[itemize,1]{label={\color{partcolor}\textbf{\textbullet}}}
\setlist[itemize,2]{label={\color{partcolor!70}\textbf{\textendash}}}
\setlist[enumerate,1]{label={\color{partcolor}\textbf{\arabic*.}}}

% =========================================================
%   PULL QUOTE / ACCENT
% =========================================================
\newcommand{\pullquote}[1]{%
  \begin{tcolorbox}[
    enhanced,
    colback=partcolor!8,
    colframe=partcolor,
    boxrule=0pt,
    leftrule=4pt,
    sharp corners,
    arc=0pt,
    fontupper=\itshape\large,
    left=14pt, right=14pt, top=8pt, bottom=8pt,
    overlay={%
      \node[partcolor!40, font=\fontsize{40}{40}\selectfont, anchor=north west]
        at ([xshift=2pt,yshift=2pt]frame.north west) {\textbf{``}};
    },
  ]
    #1
  \end{tcolorbox}
}

% =========================================================
%   DIVIDER DECORATIV (foloseste \partdivider in document)
%   Atentie: NU folosi numele "\sectionbreak" -- e rezervat
%   intern de titlesec si va aparea automat intre sectiuni.
% =========================================================
\newcommand{\partdivider}{%
  \par\vspace{0.6em}
  \begin{center}
    \tikz[baseline=-0.5ex]{
      \fill[partcolor] (0,0) circle (2pt);
      \fill[partcolor!60] (10pt,0) circle (2pt);
      \fill[partcolor!30] (20pt,0) circle (2pt);
    }
  \end{center}
  \vspace{0.3em}
}

% =========================================================
%   BOX MODEL DIAGRAM (TikZ)
% =========================================================
\newcommand{\boxmodelfigure}{%
\begin{center}
\begin{tikzpicture}[scale=1, transform shape, every node/.style={font=\sffamily\small}]
  % Margin (outermost)
  \fill[orange!25, rounded corners=4pt] (0,0) rectangle (8,5);
  \node[anchor=north west] at (0.15, 4.85) {\textbf{\color{orange!70!black}\faExpand\ margin}};

  % Border
  \fill[gray!40, rounded corners=3pt] (1,0.7) rectangle (7, 4.3);
  \draw[ultra thick, gray!70!black, rounded corners=3pt] (1,0.7) rectangle (7,4.3);
  \node[anchor=north west] at (1.1, 4.2) {\textbf{\color{gray!50!black}\faSquare\ border}};

  % Padding
  \fill[green!25, rounded corners=2pt] (1.3, 1) rectangle (6.7, 4);
  \node[anchor=north west] at (1.4, 3.9) {\textbf{\color{green!50!black}\faCompress\ padding}};

  % Content (innermost)
  \fill[blue!25, rounded corners=2pt] (2.2, 1.7) rectangle (5.8, 3.3);
  \node[anchor=north west] at (2.3, 3.2) {\textbf{\color{blue!60!black}\faFile*\ content}};
  \node at (4, 2.5) {\textit{textul, imaginea, ce e inauntru}};
\end{tikzpicture}
\end{center}
}

% =========================================================
%   CARD "ICON + TEXT" (utilitate pentru handout-uri)
% =========================================================
\newcommand{\iconcard}[3]{%
  \begin{tcolorbox}[
    enhanced,
    colback=partcolor!8,
    colframe=partcolor!30,
    boxrule=0.5pt,
    arc=6pt,
    left=12pt, right=12pt, top=8pt, bottom=8pt,
  ]
  \begin{minipage}[t]{0.08\linewidth}
    \color{partcolor}\fontsize{20}{20}\selectfont #1
  \end{minipage}%
  \begin{minipage}[t]{0.9\linewidth}
    \textbf{\color{partcolor!85!black}#2}\\
    #3
  \end{minipage}
  \end{tcolorbox}%
}
 dupa \documentclass.

\usepackage[utf8]{inputenc}
\usepackage[T1]{fontenc}
% Body font: Lato -- modern, lizibil, profesional. Fallback la TeX Gyre Heros
% (Helvetica-like, mereu disponibil) daca Lato nu e instalat in MiKTeX.
\IfFileExists{lato.sty}{%
  \usepackage[default]{lato}%
}{%
  \usepackage{tgheros}%
  \renewcommand{\familydefault}{\sfdefault}%
}
\usepackage[romanian]{babel}
\usepackage[a4paper, margin=2.5cm, headheight=16pt, top=2.4cm, bottom=2.4cm]{geometry}
\usepackage[final, protrusion=true, expansion=false]{microtype}
\usepackage{setspace}
\usepackage{xcolor}
\usepackage{amssymb}
\usepackage{listings}
\usepackage{hyperref}
\usepackage{enumitem}
\usepackage{titlesec}
\usepackage{fancyhdr}
\usepackage{parskip}
\usepackage[most]{tcolorbox}
\usepackage{tocloft}
\usepackage{graphicx}
\usepackage{tikz}
\usepackage{fontawesome5}
\usepackage{lettrine}
\usetikzlibrary{shadows.blur, shapes.geometric, positioning, calc,
  decorations.pathmorphing, patterns, fadings, backgrounds}

% =========================================================
%   PALETA DE CULORI
% =========================================================

% Paleta finala (3 BEST + 3 derivate in acelasi ton)
\definecolor{themePrereq}{HTML}{5C6B7A}    % slate (pre-rechizite, neutral)
\definecolor{themeHtml}{HTML}{FAA51E}      % BEST orange
\definecolor{themeCss}{HTML}{0673BA}       % BEST blue
\definecolor{themeJs}{HTML}{ECC100}        % gold (analog cu orange, distinct)
\definecolor{themeDeploy}{HTML}{75C044}    % BEST green
\definecolor{themeSeo}{HTML}{B344C9}       % violet (analog cu blue)
\definecolor{themeLighthouse}{HTML}{DD4444}% coral red (complementar cu green)
\definecolor{themeResurse}{HTML}{5C6B7A}   % slate (resurse, neutral)
\definecolor{themeBonus}{HTML}{7E22CE}     % deep purple (bonus / avansat, self-study only)
% Aliasuri pentru compatibilitate (folosit in main.tex existent)
\colorlet{themeIntro}{themeHtml}
\colorlet{themeExercitiu}{themeLighthouse}

% Background pagina si elemente -- calduros
\definecolor{pageBg}{HTML}{FFFCF6}         % crema foarte subtila
\definecolor{accentSoft}{HTML}{F4ECDD}     % accent cald
\definecolor{textDim}{HTML}{64748B}        % gri pentru text secundar

% Culoarea curenta a partii (suprascrisa de fiecare fisier prin \setpart)
\colorlet{partcolor}{themeIntro}

% Iconita partii curente (suprascrisa prin \setpart)
\newcommand{\particon}{\faStar}
\newcommand{\partlabel}{Training Web}

% \setpart{culoare}{iconita}
\newcommand{\setpart}[2]{%
  \colorlet{partcolor}{#1}%
  \renewcommand{\particon}{#2}%
}
% \setpart{culoare}{iconita}{eticheta} -- versiunea cu eticheta custom in header
\newcommand{\setpartlabel}[3]{%
  \colorlet{partcolor}{#1}%
  \renewcommand{\particon}{#2}%
  \renewcommand{\partlabel}{#3}%
}

% Culori pentru callout boxes
\definecolor{noteBg}{HTML}{FEF9C3}        \definecolor{noteFg}{HTML}{A16207}
\definecolor{tipBg}{HTML}{DBEAFE}         \definecolor{tipFg}{HTML}{1D4ED8}
\definecolor{warnBg}{HTML}{FEE2E2}        \definecolor{warnFg}{HTML}{B91C1C}
\definecolor{gameBg}{HTML}{FFEDD5}        \definecolor{gameFg}{HTML}{C2410C}
\definecolor{recapBg}{HTML}{DCFCE7}       \definecolor{recapFg}{HTML}{15803D}
\definecolor{exerBg}{HTML}{F3E8FF}        \definecolor{exerFg}{HTML}{7C3AED}
\definecolor{triviaBg}{HTML}{FCE7F3}      \definecolor{triviaFg}{HTML}{BE185D}

% Culori pentru cod -- paleta "GitHub Light" (calm, profesional)
\definecolor{codeBg}{HTML}{F6F8FA}
\definecolor{codeBorder}{HTML}{D0D7DE}
\definecolor{codeComment}{HTML}{6E7781}
\definecolor{codeKeyword}{HTML}{CF222E}
\definecolor{codeString}{HTML}{0A3069}
\definecolor{codeTag}{HTML}{116329}

% =========================================================
%   PAGINA & LIZIBILITATE (warmth)
% =========================================================
\pagecolor{pageBg}
\setstretch{1.18}
\setlength{\parskip}{0.5em}

% =========================================================
%   CUPRINS (TOC)
% =========================================================
\setcounter{tocdepth}{0}         % numai entry-uri \part (capitole), fara sectiuni
\renewcommand{\contentsname}{Cuprins}

% Stil cuprins -- numai capitolele, asezate elegant
\renewcommand{\cftpartfont}{\sffamily\bfseries\color{themeText}}
\renewcommand{\cftpartpagefont}{\sffamily\bfseries\color{themeText}}
\renewcommand{\cftpartleader}{\hfill}
\renewcommand{\cftpartafterpnum}{\par\addvspace{2pt}}
\setlength{\cftbeforepartskip}{4pt}

% =========================================================
%   HYPERREF
% =========================================================
\hypersetup{
  colorlinks=true,
  linkcolor=partcolor,
  urlcolor=partcolor,
  citecolor=partcolor,
  pdfborderstyle={/S/U/W 1}
}

% =========================================================
%   HEADER & FOOTER -- minimaliste
% =========================================================
\pagestyle{fancy}
\fancyhf{}
\fancyhead[L]{\footnotesize\sffamily\color{black!55}\partlabel}
\fancyhead[R]{\footnotesize\sffamily\color{black!55}BEST Ia\textcommabelow{s}i}
\fancyfoot[C]{\small\sffamily\color{partcolor}\textbf{\thepage}}
\renewcommand{\headrulewidth}{0pt}
\renewcommand{\footrulewidth}{0pt}

% =========================================================
%   STIL TITLURI -- sobru, stiintific
% =========================================================
\titleformat{\section}
  {\Large\sffamily\bfseries\color{themeText}}
  {\color{partcolor}\sffamily\bfseries\thesection.\hspace{0.5em}}
  {0pt}{}
\titlespacing*{\section}{0pt}{1.4em}{0.4em}

\titleformat{\subsection}
  {\normalsize\sffamily\bfseries\color{black!75}}
  {\color{partcolor}\thesubsection\hspace{0.5em}}{0pt}{}
\titlespacing*{\subsection}{0pt}{0.8em}{0.2em}

\titleformat{\subsubsection}
  {\normalsize\sffamily\bfseries\color{black!70}}
  {}{0pt}{}

\definecolor{themeText}{HTML}{1A1A1A}

% =========================================================
%   LISTINGS -- COD
% =========================================================
\lstdefinestyle{base}{
  backgroundcolor=\color{codeBg},
  basicstyle=\ttfamily\small,
  commentstyle=\color{codeComment}\itshape,
  keywordstyle=\color{codeKeyword}\bfseries,
  stringstyle=\color{codeString},
  numberstyle=\tiny\color{codeComment},
  numbers=left,
  numbersep=8pt,
  frame=leftline,
  framerule=2.5pt,
  rulecolor=\color{partcolor},
  framexleftmargin=4pt,
  breaklines=true,
  breakatwhitespace=true,
  showstringspaces=false,
  tabsize=2,
  xleftmargin=14pt,
  xrightmargin=4pt,
  aboveskip=10pt,
  belowskip=10pt,
  literate=
    {ă}{{\u a}}1 {Ă}{{\u A}}1
    {â}{{\^a}}1 {Â}{{\^A}}1
    {î}{{\^\i}}1 {Î}{{\^I}}1
    {ș}{{\textcommabelow s}}1 {Ș}{{\textcommabelow S}}1
    {ț}{{\textcommabelow t}}1 {Ț}{{\textcommabelow T}}1
}

\lstdefinelanguage{HTML5}{
  morekeywords={html,head,body,title,meta,link,script,style,
    h1,h2,h3,h4,h5,h6,p,a,img,div,span,
    header,main,section,article,aside,nav,footer,
    ul,ol,li,button,input,form,label,
    strong,em,br,hr,DOCTYPE},
  sensitive=true,
  morecomment=[s]{<!--}{-->},
  morestring=[b]",
  morestring=[b]'
}
\lstdefinestyle{html}{style=base, language=HTML5,
  keywordstyle=\color{codeTag}\bfseries}

\lstdefinelanguage{CSSLang}{
  morekeywords={display,flex,grid,block,inline,none,
    color,background,background-color,
    padding,padding-top,padding-right,padding-bottom,padding-left,
    margin,margin-top,margin-right,margin-bottom,margin-left,
    border,border-radius,border-color,outline,
    width,max-width,min-width,height,max-height,min-height,
    font-family,font-size,font-weight,line-height,text-align,
    text-decoration,letter-spacing,
    flex-direction,justify-content,align-items,gap,flex-wrap,flex-grow,
    grid-template-columns,grid-template-rows,grid-column,grid-row,
    position,top,right,bottom,left,z-index,
    box-sizing,opacity,transition,transform,
    cursor,overflow,visibility,
    media,keyframes,from,to,important,root,hover,focus,active,
    auto,solid,dashed,dotted,bold,normal,center,stretch,
    space-between,space-around,flex-start,flex-end,column,row,wrap},
  sensitive=true,
  morecomment=[s]{/*}{*/},
  morestring=[b]",
  morestring=[b]'
}
\lstdefinestyle{css}{style=base, language=CSSLang,
  keywordstyle=\color{codeKeyword}\bfseries}

\lstdefinelanguage{JSLang}{
  morekeywords={let,const,var,function,return,if,else,for,while,
    true,false,null,undefined,new,this,class,extends,
    document,console,window,
    addEventListener,querySelector,querySelectorAll,
    getElementById,getElementsByClassName,
    classList,toggle,add,remove,contains,
    style,textContent,innerHTML,value,
    log,warn,error},
  sensitive=true,
  morecomment=[l]{//},
  morecomment=[s]{/*}{*/},
  morestring=[b]",
  morestring=[b]',
  morestring=[b]`
}
\lstdefinestyle{js}{style=base, language=JSLang,
  keywordstyle=\color{codeKeyword}\bfseries}

\lstdefinelanguage{BashLang}{
  morekeywords={git,cd,mkdir,ls,code,echo,cat,curl,wget,sudo,
    init,add,commit,push,pull,clone,remote,branch,checkout,merge,status,log,
    config,fetch,rebase,reset,stash,diff,tag,
    install,update,run,build,start,test},
  sensitive=true,
  morecomment=[l]{\#},
  morestring=[b]",
  morestring=[b]'
}
\lstdefinestyle{bash}{style=base, language=BashLang,
  keywordstyle=\color{codeKeyword}\bfseries}

% Inline code -- mai discret
\newcommand{\code}[1]{\textcolor{partcolor!85!black}{\ttfamily\small #1}}

% =========================================================
%   CALLOUT BOXES -- minimaliste (bara stanga subtila + label discret)
% =========================================================
\tcbset{
  calloutbase/.style={
    enhanced,
    breakable,
    sharp corners,
    boxrule=0pt,
    leftrule=2.5pt,
    arc=0pt,
    colback=pageBg,
    fonttitle=\sffamily\small\bfseries,
    coltitle=tcbcolframe,
    attach title to upper={\par\vspace{2pt}},
    title after break={},
    before skip=10pt,
    after skip=10pt,
    left=12pt, right=4pt, top=4pt, bottom=4pt,
    fontupper=\small,
  }
}

\newtcolorbox{nota}[1][NOT\u A]{
  calloutbase,
  colframe=noteFg,
  title={\faInfoCircle\ \uppercase\expandafter{#1}},
}
\newtcolorbox{ponturi}[1][PONT]{
  calloutbase,
  colframe=tipFg,
  title={\faLightbulb[regular]\ \uppercase\expandafter{#1}},
}
\newtcolorbox{atentie}[1][ATEN\textcommabelow{T}IE]{
  calloutbase,
  colframe=warnFg,
  title={\faExclamationTriangle\ \uppercase\expandafter{#1}},
}
\newtcolorbox{joc}[1][APLICA\textcommabelow{T}IE PRACTIC\u A]{
  calloutbase,
  colframe=partcolor,
  title={\faGamepad\ \uppercase\expandafter{#1}},
}
\newtcolorbox{recap}[1][RECAP]{
  calloutbase,
  colframe=recapFg,
  title={\faCheckCircle\ \uppercase\expandafter{#1}},
}
\newtcolorbox{curiozitate}[1][\textcommabelow{S}TIAI C\u A?]{
  calloutbase,
  colframe=triviaFg,
  title={\faMagic\ \uppercase\expandafter{#1}},
}
\newtcolorbox{exercitiu}[1][EXERCI\textcommabelow{T}IU]{
  calloutbase,
  colframe=exerFg,
  title={\faPen\ \uppercase\expandafter{#1}},
}

% --- Outcome / Why / Self-check (structura tutorial) ----
\definecolor{outBg}{HTML}{E0F2FE}        \definecolor{outFg}{HTML}{0369A1}
\definecolor{whyBg}{HTML}{FEF3C7}        \definecolor{whyFg}{HTML}{B45309}
\definecolor{checkBg}{HTML}{F3E8FF}      \definecolor{checkFg}{HTML}{7C3AED}

\newtcolorbox{outcomes}[1][CE INVE\textcommabelow{T}I]{
  calloutbase,
  colframe=outFg,
  title={\faBullseye\ \uppercase\expandafter{#1}},
}
\newtcolorbox{whymatters}[1][DE CE]{
  calloutbase,
  colframe=whyFg,
  title={\faQuestionCircle\ \uppercase\expandafter{#1}},
}
\newtcolorbox{selfcheck}[1][PROVOCARE]{
  calloutbase,
  colframe=checkFg,
  title={\faTasks\ \uppercase\expandafter{#1}},
}

% --- Bonus / Avansat (continut auto-studiu, distinct vizual) ----
% Diferit fata de celelalte callout-uri: fundal usor colorat (nu pageBg),
% bara stanga mai groasa, icon racheta. Beginners vad clar si pot sari peste.
\definecolor{bonusBg}{HTML}{F6EEFB}        \definecolor{bonusFg}{HTML}{6B21A8}

\newtcolorbox{bonusbox}[1][BONUS \,\textbullet\, AVANSAT]{
  enhanced,
  breakable,
  sharp corners,
  boxrule=0pt,
  leftrule=4pt,
  arc=0pt,
  colback=bonusBg,
  colframe=themeBonus,
  fonttitle=\sffamily\small\bfseries,
  coltitle=bonusFg,
  attach title to upper={\par\vspace{2pt}},
  title after break={\faRocket\ \textbf{\color{bonusFg}continuare bonus}},
  before skip=14pt,
  after skip=14pt,
  left=14pt, right=6pt, top=6pt, bottom=6pt,
  fontupper=\small,
  title={\faRocket\ \uppercase\expandafter{#1}},
}

% Bridge la finalul fiecarui capitol
\newcommand{\bridge}[1]{%
  \par\vspace{4pt}%
  \begin{flushleft}%
    {\sffamily\small\color{partcolor}$\rightarrow$\ \textbf{Urmeaza:}}\quad{\small\color{black!70}#1}%
  \end{flushleft}%
}

% --- Asa da / Asa nu (compare bune vs rele practici) ----
\definecolor{daBg}{HTML}{DCFCE7}        \definecolor{daFg}{HTML}{15803D}
\definecolor{nuBg}{HTML}{FEE2E2}        \definecolor{nuFg}{HTML}{B91C1C}

\newtcolorbox{dado}[1][A\textcommabelow{S}A DA]{
  calloutbase,
  colframe=daFg,
  title={\faCheckCircle\ \uppercase\expandafter{#1}},
}
\newtcolorbox{dano}[1][A\textcommabelow{S}A NU]{
  calloutbase,
  colframe=nuFg,
  title={\faTimesCircle\ \uppercase\expandafter{#1}},
}

% =========================================================
%   BANNER MINIMALIST (\handout{titlu}{subtitlu})
%   Eticheta sus, titlu mare, subtitlu, linie subtila colorata.
% =========================================================
\newcommand{\handout}[2]{%
  \noindent\begin{minipage}{0.04\linewidth}\color{partcolor}\rule{4pt}{2.6em}\end{minipage}%
  \begin{minipage}{0.94\linewidth}
    {\sffamily\fontsize{30}{34}\selectfont\textbf{\color{themeText}#1}}\\[4pt]
    {\sffamily\large\color{black!55}#2}
  \end{minipage}
  \par\vspace{14pt}
  {\color{partcolor!40}\hrule height 0.5pt}
  \par\vspace{14pt}
}

% =========================================================
%   BULLETS COLORATI (itemize)
% =========================================================
\setlist[itemize,1]{label={\color{partcolor}\textbf{\textbullet}}}
\setlist[itemize,2]{label={\color{partcolor!70}\textbf{\textendash}}}
\setlist[enumerate,1]{label={\color{partcolor}\textbf{\arabic*.}}}

% =========================================================
%   PULL QUOTE / ACCENT
% =========================================================
\newcommand{\pullquote}[1]{%
  \begin{tcolorbox}[
    enhanced,
    colback=partcolor!8,
    colframe=partcolor,
    boxrule=0pt,
    leftrule=4pt,
    sharp corners,
    arc=0pt,
    fontupper=\itshape\large,
    left=14pt, right=14pt, top=8pt, bottom=8pt,
    overlay={%
      \node[partcolor!40, font=\fontsize{40}{40}\selectfont, anchor=north west]
        at ([xshift=2pt,yshift=2pt]frame.north west) {\textbf{``}};
    },
  ]
    #1
  \end{tcolorbox}
}

% =========================================================
%   DIVIDER DECORATIV (foloseste \partdivider in document)
%   Atentie: NU folosi numele "\sectionbreak" -- e rezervat
%   intern de titlesec si va aparea automat intre sectiuni.
% =========================================================
\newcommand{\partdivider}{%
  \par\vspace{0.6em}
  \begin{center}
    \tikz[baseline=-0.5ex]{
      \fill[partcolor] (0,0) circle (2pt);
      \fill[partcolor!60] (10pt,0) circle (2pt);
      \fill[partcolor!30] (20pt,0) circle (2pt);
    }
  \end{center}
  \vspace{0.3em}
}

% =========================================================
%   BOX MODEL DIAGRAM (TikZ)
% =========================================================
\newcommand{\boxmodelfigure}{%
\begin{center}
\begin{tikzpicture}[scale=1, transform shape, every node/.style={font=\sffamily\small}]
  % Margin (outermost)
  \fill[orange!25, rounded corners=4pt] (0,0) rectangle (8,5);
  \node[anchor=north west] at (0.15, 4.85) {\textbf{\color{orange!70!black}\faExpand\ margin}};

  % Border
  \fill[gray!40, rounded corners=3pt] (1,0.7) rectangle (7, 4.3);
  \draw[ultra thick, gray!70!black, rounded corners=3pt] (1,0.7) rectangle (7,4.3);
  \node[anchor=north west] at (1.1, 4.2) {\textbf{\color{gray!50!black}\faSquare\ border}};

  % Padding
  \fill[green!25, rounded corners=2pt] (1.3, 1) rectangle (6.7, 4);
  \node[anchor=north west] at (1.4, 3.9) {\textbf{\color{green!50!black}\faCompress\ padding}};

  % Content (innermost)
  \fill[blue!25, rounded corners=2pt] (2.2, 1.7) rectangle (5.8, 3.3);
  \node[anchor=north west] at (2.3, 3.2) {\textbf{\color{blue!60!black}\faFile*\ content}};
  \node at (4, 2.5) {\textit{textul, imaginea, ce e inauntru}};
\end{tikzpicture}
\end{center}
}

% =========================================================
%   CARD "ICON + TEXT" (utilitate pentru handout-uri)
% =========================================================
\newcommand{\iconcard}[3]{%
  \begin{tcolorbox}[
    enhanced,
    colback=partcolor!8,
    colframe=partcolor!30,
    boxrule=0.5pt,
    arc=6pt,
    left=12pt, right=12pt, top=8pt, bottom=8pt,
  ]
  \begin{minipage}[t]{0.08\linewidth}
    \color{partcolor}\fontsize{20}{20}\selectfont #1
  \end{minipage}%
  \begin{minipage}[t]{0.9\linewidth}
    \textbf{\color{partcolor!85!black}#2}\\
    #3
  \end{minipage}
  \end{tcolorbox}%
}
 dupa \documentclass.

\usepackage[utf8]{inputenc}
\usepackage[T1]{fontenc}
% Body font: Lato -- modern, lizibil, profesional. Fallback la TeX Gyre Heros
% (Helvetica-like, mereu disponibil) daca Lato nu e instalat in MiKTeX.
\IfFileExists{lato.sty}{%
  \usepackage[default]{lato}%
}{%
  \usepackage{tgheros}%
  \renewcommand{\familydefault}{\sfdefault}%
}
\usepackage[romanian]{babel}
\usepackage[a4paper, margin=2.5cm, headheight=16pt, top=2.4cm, bottom=2.4cm]{geometry}
\usepackage[final, protrusion=true, expansion=false]{microtype}
\usepackage{setspace}
\usepackage{xcolor}
\usepackage{amssymb}
\usepackage{listings}
\usepackage{hyperref}
\usepackage{enumitem}
\usepackage{titlesec}
\usepackage{fancyhdr}
\usepackage{parskip}
\usepackage[most]{tcolorbox}
\usepackage{tocloft}
\usepackage{graphicx}
\usepackage{tikz}
\usepackage{fontawesome5}
\usepackage{lettrine}
\usetikzlibrary{shadows.blur, shapes.geometric, positioning, calc,
  decorations.pathmorphing, patterns, fadings, backgrounds}

% =========================================================
%   PALETA DE CULORI
% =========================================================

% Paleta finala (3 BEST + 3 derivate in acelasi ton)
\definecolor{themePrereq}{HTML}{5C6B7A}    % slate (pre-rechizite, neutral)
\definecolor{themeHtml}{HTML}{FAA51E}      % BEST orange
\definecolor{themeCss}{HTML}{0673BA}       % BEST blue
\definecolor{themeJs}{HTML}{ECC100}        % gold (analog cu orange, distinct)
\definecolor{themeDeploy}{HTML}{75C044}    % BEST green
\definecolor{themeSeo}{HTML}{B344C9}       % violet (analog cu blue)
\definecolor{themeLighthouse}{HTML}{DD4444}% coral red (complementar cu green)
\definecolor{themeResurse}{HTML}{5C6B7A}   % slate (resurse, neutral)
\definecolor{themeBonus}{HTML}{7E22CE}     % deep purple (bonus / avansat, self-study only)
% Aliasuri pentru compatibilitate (folosit in main.tex existent)
\colorlet{themeIntro}{themeHtml}
\colorlet{themeExercitiu}{themeLighthouse}

% Background pagina si elemente -- calduros
\definecolor{pageBg}{HTML}{FFFCF6}         % crema foarte subtila
\definecolor{accentSoft}{HTML}{F4ECDD}     % accent cald
\definecolor{textDim}{HTML}{64748B}        % gri pentru text secundar

% Culoarea curenta a partii (suprascrisa de fiecare fisier prin \setpart)
\colorlet{partcolor}{themeIntro}

% Iconita partii curente (suprascrisa prin \setpart)
\newcommand{\particon}{\faStar}
\newcommand{\partlabel}{Training Web}

% \setpart{culoare}{iconita}
\newcommand{\setpart}[2]{%
  \colorlet{partcolor}{#1}%
  \renewcommand{\particon}{#2}%
}
% \setpart{culoare}{iconita}{eticheta} -- versiunea cu eticheta custom in header
\newcommand{\setpartlabel}[3]{%
  \colorlet{partcolor}{#1}%
  \renewcommand{\particon}{#2}%
  \renewcommand{\partlabel}{#3}%
}

% Culori pentru callout boxes
\definecolor{noteBg}{HTML}{FEF9C3}        \definecolor{noteFg}{HTML}{A16207}
\definecolor{tipBg}{HTML}{DBEAFE}         \definecolor{tipFg}{HTML}{1D4ED8}
\definecolor{warnBg}{HTML}{FEE2E2}        \definecolor{warnFg}{HTML}{B91C1C}
\definecolor{gameBg}{HTML}{FFEDD5}        \definecolor{gameFg}{HTML}{C2410C}
\definecolor{recapBg}{HTML}{DCFCE7}       \definecolor{recapFg}{HTML}{15803D}
\definecolor{exerBg}{HTML}{F3E8FF}        \definecolor{exerFg}{HTML}{7C3AED}
\definecolor{triviaBg}{HTML}{FCE7F3}      \definecolor{triviaFg}{HTML}{BE185D}

% Culori pentru cod -- paleta "GitHub Light" (calm, profesional)
\definecolor{codeBg}{HTML}{F6F8FA}
\definecolor{codeBorder}{HTML}{D0D7DE}
\definecolor{codeComment}{HTML}{6E7781}
\definecolor{codeKeyword}{HTML}{CF222E}
\definecolor{codeString}{HTML}{0A3069}
\definecolor{codeTag}{HTML}{116329}

% =========================================================
%   PAGINA & LIZIBILITATE (warmth)
% =========================================================
\pagecolor{pageBg}
\setstretch{1.18}
\setlength{\parskip}{0.5em}

% =========================================================
%   CUPRINS (TOC)
% =========================================================
\setcounter{tocdepth}{0}         % numai entry-uri \part (capitole), fara sectiuni
\renewcommand{\contentsname}{Cuprins}

% Stil cuprins -- numai capitolele, asezate elegant
\renewcommand{\cftpartfont}{\sffamily\bfseries\color{themeText}}
\renewcommand{\cftpartpagefont}{\sffamily\bfseries\color{themeText}}
\renewcommand{\cftpartleader}{\hfill}
\renewcommand{\cftpartafterpnum}{\par\addvspace{2pt}}
\setlength{\cftbeforepartskip}{4pt}

% =========================================================
%   HYPERREF
% =========================================================
\hypersetup{
  colorlinks=true,
  linkcolor=partcolor,
  urlcolor=partcolor,
  citecolor=partcolor,
  pdfborderstyle={/S/U/W 1}
}

% =========================================================
%   HEADER & FOOTER -- minimaliste
% =========================================================
\pagestyle{fancy}
\fancyhf{}
\fancyhead[L]{\footnotesize\sffamily\color{black!55}\partlabel}
\fancyhead[R]{\footnotesize\sffamily\color{black!55}BEST Ia\textcommabelow{s}i}
\fancyfoot[C]{\small\sffamily\color{partcolor}\textbf{\thepage}}
\renewcommand{\headrulewidth}{0pt}
\renewcommand{\footrulewidth}{0pt}

% =========================================================
%   STIL TITLURI -- sobru, stiintific
% =========================================================
\titleformat{\section}
  {\Large\sffamily\bfseries\color{themeText}}
  {\color{partcolor}\sffamily\bfseries\thesection.\hspace{0.5em}}
  {0pt}{}
\titlespacing*{\section}{0pt}{1.4em}{0.4em}

\titleformat{\subsection}
  {\normalsize\sffamily\bfseries\color{black!75}}
  {\color{partcolor}\thesubsection\hspace{0.5em}}{0pt}{}
\titlespacing*{\subsection}{0pt}{0.8em}{0.2em}

\titleformat{\subsubsection}
  {\normalsize\sffamily\bfseries\color{black!70}}
  {}{0pt}{}

\definecolor{themeText}{HTML}{1A1A1A}

% =========================================================
%   LISTINGS -- COD
% =========================================================
\lstdefinestyle{base}{
  backgroundcolor=\color{codeBg},
  basicstyle=\ttfamily\small,
  commentstyle=\color{codeComment}\itshape,
  keywordstyle=\color{codeKeyword}\bfseries,
  stringstyle=\color{codeString},
  numberstyle=\tiny\color{codeComment},
  numbers=left,
  numbersep=8pt,
  frame=leftline,
  framerule=2.5pt,
  rulecolor=\color{partcolor},
  framexleftmargin=4pt,
  breaklines=true,
  breakatwhitespace=true,
  showstringspaces=false,
  tabsize=2,
  xleftmargin=14pt,
  xrightmargin=4pt,
  aboveskip=10pt,
  belowskip=10pt,
  literate=
    {ă}{{\u a}}1 {Ă}{{\u A}}1
    {â}{{\^a}}1 {Â}{{\^A}}1
    {î}{{\^\i}}1 {Î}{{\^I}}1
    {ș}{{\textcommabelow s}}1 {Ș}{{\textcommabelow S}}1
    {ț}{{\textcommabelow t}}1 {Ț}{{\textcommabelow T}}1
}

\lstdefinelanguage{HTML5}{
  morekeywords={html,head,body,title,meta,link,script,style,
    h1,h2,h3,h4,h5,h6,p,a,img,div,span,
    header,main,section,article,aside,nav,footer,
    ul,ol,li,button,input,form,label,
    strong,em,br,hr,DOCTYPE},
  sensitive=true,
  morecomment=[s]{<!--}{-->},
  morestring=[b]",
  morestring=[b]'
}
\lstdefinestyle{html}{style=base, language=HTML5,
  keywordstyle=\color{codeTag}\bfseries}

\lstdefinelanguage{CSSLang}{
  morekeywords={display,flex,grid,block,inline,none,
    color,background,background-color,
    padding,padding-top,padding-right,padding-bottom,padding-left,
    margin,margin-top,margin-right,margin-bottom,margin-left,
    border,border-radius,border-color,outline,
    width,max-width,min-width,height,max-height,min-height,
    font-family,font-size,font-weight,line-height,text-align,
    text-decoration,letter-spacing,
    flex-direction,justify-content,align-items,gap,flex-wrap,flex-grow,
    grid-template-columns,grid-template-rows,grid-column,grid-row,
    position,top,right,bottom,left,z-index,
    box-sizing,opacity,transition,transform,
    cursor,overflow,visibility,
    media,keyframes,from,to,important,root,hover,focus,active,
    auto,solid,dashed,dotted,bold,normal,center,stretch,
    space-between,space-around,flex-start,flex-end,column,row,wrap},
  sensitive=true,
  morecomment=[s]{/*}{*/},
  morestring=[b]",
  morestring=[b]'
}
\lstdefinestyle{css}{style=base, language=CSSLang,
  keywordstyle=\color{codeKeyword}\bfseries}

\lstdefinelanguage{JSLang}{
  morekeywords={let,const,var,function,return,if,else,for,while,
    true,false,null,undefined,new,this,class,extends,
    document,console,window,
    addEventListener,querySelector,querySelectorAll,
    getElementById,getElementsByClassName,
    classList,toggle,add,remove,contains,
    style,textContent,innerHTML,value,
    log,warn,error},
  sensitive=true,
  morecomment=[l]{//},
  morecomment=[s]{/*}{*/},
  morestring=[b]",
  morestring=[b]',
  morestring=[b]`
}
\lstdefinestyle{js}{style=base, language=JSLang,
  keywordstyle=\color{codeKeyword}\bfseries}

\lstdefinelanguage{BashLang}{
  morekeywords={git,cd,mkdir,ls,code,echo,cat,curl,wget,sudo,
    init,add,commit,push,pull,clone,remote,branch,checkout,merge,status,log,
    config,fetch,rebase,reset,stash,diff,tag,
    install,update,run,build,start,test},
  sensitive=true,
  morecomment=[l]{\#},
  morestring=[b]",
  morestring=[b]'
}
\lstdefinestyle{bash}{style=base, language=BashLang,
  keywordstyle=\color{codeKeyword}\bfseries}

% Inline code -- mai discret
\newcommand{\code}[1]{\textcolor{partcolor!85!black}{\ttfamily\small #1}}

% =========================================================
%   CALLOUT BOXES -- minimaliste (bara stanga subtila + label discret)
% =========================================================
\tcbset{
  calloutbase/.style={
    enhanced,
    breakable,
    sharp corners,
    boxrule=0pt,
    leftrule=2.5pt,
    arc=0pt,
    colback=pageBg,
    fonttitle=\sffamily\small\bfseries,
    coltitle=tcbcolframe,
    attach title to upper={\par\vspace{2pt}},
    title after break={},
    before skip=10pt,
    after skip=10pt,
    left=12pt, right=4pt, top=4pt, bottom=4pt,
    fontupper=\small,
  }
}

\newtcolorbox{nota}[1][NOT\u A]{
  calloutbase,
  colframe=noteFg,
  title={\faInfoCircle\ \uppercase\expandafter{#1}},
}
\newtcolorbox{ponturi}[1][PONT]{
  calloutbase,
  colframe=tipFg,
  title={\faLightbulb[regular]\ \uppercase\expandafter{#1}},
}
\newtcolorbox{atentie}[1][ATEN\textcommabelow{T}IE]{
  calloutbase,
  colframe=warnFg,
  title={\faExclamationTriangle\ \uppercase\expandafter{#1}},
}
\newtcolorbox{joc}[1][APLICA\textcommabelow{T}IE PRACTIC\u A]{
  calloutbase,
  colframe=partcolor,
  title={\faGamepad\ \uppercase\expandafter{#1}},
}
\newtcolorbox{recap}[1][RECAP]{
  calloutbase,
  colframe=recapFg,
  title={\faCheckCircle\ \uppercase\expandafter{#1}},
}
\newtcolorbox{curiozitate}[1][\textcommabelow{S}TIAI C\u A?]{
  calloutbase,
  colframe=triviaFg,
  title={\faMagic\ \uppercase\expandafter{#1}},
}
\newtcolorbox{exercitiu}[1][EXERCI\textcommabelow{T}IU]{
  calloutbase,
  colframe=exerFg,
  title={\faPen\ \uppercase\expandafter{#1}},
}

% --- Outcome / Why / Self-check (structura tutorial) ----
\definecolor{outBg}{HTML}{E0F2FE}        \definecolor{outFg}{HTML}{0369A1}
\definecolor{whyBg}{HTML}{FEF3C7}        \definecolor{whyFg}{HTML}{B45309}
\definecolor{checkBg}{HTML}{F3E8FF}      \definecolor{checkFg}{HTML}{7C3AED}

\newtcolorbox{outcomes}[1][CE INVE\textcommabelow{T}I]{
  calloutbase,
  colframe=outFg,
  title={\faBullseye\ \uppercase\expandafter{#1}},
}
\newtcolorbox{whymatters}[1][DE CE]{
  calloutbase,
  colframe=whyFg,
  title={\faQuestionCircle\ \uppercase\expandafter{#1}},
}
\newtcolorbox{selfcheck}[1][PROVOCARE]{
  calloutbase,
  colframe=checkFg,
  title={\faTasks\ \uppercase\expandafter{#1}},
}

% --- Bonus / Avansat (continut auto-studiu, distinct vizual) ----
% Diferit fata de celelalte callout-uri: fundal usor colorat (nu pageBg),
% bara stanga mai groasa, icon racheta. Beginners vad clar si pot sari peste.
\definecolor{bonusBg}{HTML}{F6EEFB}        \definecolor{bonusFg}{HTML}{6B21A8}

\newtcolorbox{bonusbox}[1][BONUS \,\textbullet\, AVANSAT]{
  enhanced,
  breakable,
  sharp corners,
  boxrule=0pt,
  leftrule=4pt,
  arc=0pt,
  colback=bonusBg,
  colframe=themeBonus,
  fonttitle=\sffamily\small\bfseries,
  coltitle=bonusFg,
  attach title to upper={\par\vspace{2pt}},
  title after break={\faRocket\ \textbf{\color{bonusFg}continuare bonus}},
  before skip=14pt,
  after skip=14pt,
  left=14pt, right=6pt, top=6pt, bottom=6pt,
  fontupper=\small,
  title={\faRocket\ \uppercase\expandafter{#1}},
}

% Bridge la finalul fiecarui capitol
\newcommand{\bridge}[1]{%
  \par\vspace{4pt}%
  \begin{flushleft}%
    {\sffamily\small\color{partcolor}$\rightarrow$\ \textbf{Urmeaza:}}\quad{\small\color{black!70}#1}%
  \end{flushleft}%
}

% --- Asa da / Asa nu (compare bune vs rele practici) ----
\definecolor{daBg}{HTML}{DCFCE7}        \definecolor{daFg}{HTML}{15803D}
\definecolor{nuBg}{HTML}{FEE2E2}        \definecolor{nuFg}{HTML}{B91C1C}

\newtcolorbox{dado}[1][A\textcommabelow{S}A DA]{
  calloutbase,
  colframe=daFg,
  title={\faCheckCircle\ \uppercase\expandafter{#1}},
}
\newtcolorbox{dano}[1][A\textcommabelow{S}A NU]{
  calloutbase,
  colframe=nuFg,
  title={\faTimesCircle\ \uppercase\expandafter{#1}},
}

% =========================================================
%   BANNER MINIMALIST (\handout{titlu}{subtitlu})
%   Eticheta sus, titlu mare, subtitlu, linie subtila colorata.
% =========================================================
\newcommand{\handout}[2]{%
  \noindent\begin{minipage}{0.04\linewidth}\color{partcolor}\rule{4pt}{2.6em}\end{minipage}%
  \begin{minipage}{0.94\linewidth}
    {\sffamily\fontsize{30}{34}\selectfont\textbf{\color{themeText}#1}}\\[4pt]
    {\sffamily\large\color{black!55}#2}
  \end{minipage}
  \par\vspace{14pt}
  {\color{partcolor!40}\hrule height 0.5pt}
  \par\vspace{14pt}
}

% =========================================================
%   BULLETS COLORATI (itemize)
% =========================================================
\setlist[itemize,1]{label={\color{partcolor}\textbf{\textbullet}}}
\setlist[itemize,2]{label={\color{partcolor!70}\textbf{\textendash}}}
\setlist[enumerate,1]{label={\color{partcolor}\textbf{\arabic*.}}}

% =========================================================
%   PULL QUOTE / ACCENT
% =========================================================
\newcommand{\pullquote}[1]{%
  \begin{tcolorbox}[
    enhanced,
    colback=partcolor!8,
    colframe=partcolor,
    boxrule=0pt,
    leftrule=4pt,
    sharp corners,
    arc=0pt,
    fontupper=\itshape\large,
    left=14pt, right=14pt, top=8pt, bottom=8pt,
    overlay={%
      \node[partcolor!40, font=\fontsize{40}{40}\selectfont, anchor=north west]
        at ([xshift=2pt,yshift=2pt]frame.north west) {\textbf{``}};
    },
  ]
    #1
  \end{tcolorbox}
}

% =========================================================
%   DIVIDER DECORATIV (foloseste \partdivider in document)
%   Atentie: NU folosi numele "\sectionbreak" -- e rezervat
%   intern de titlesec si va aparea automat intre sectiuni.
% =========================================================
\newcommand{\partdivider}{%
  \par\vspace{0.6em}
  \begin{center}
    \tikz[baseline=-0.5ex]{
      \fill[partcolor] (0,0) circle (2pt);
      \fill[partcolor!60] (10pt,0) circle (2pt);
      \fill[partcolor!30] (20pt,0) circle (2pt);
    }
  \end{center}
  \vspace{0.3em}
}

% =========================================================
%   BOX MODEL DIAGRAM (TikZ)
% =========================================================
\newcommand{\boxmodelfigure}{%
\begin{center}
\begin{tikzpicture}[scale=1, transform shape, every node/.style={font=\sffamily\small}]
  % Margin (outermost)
  \fill[orange!25, rounded corners=4pt] (0,0) rectangle (8,5);
  \node[anchor=north west] at (0.15, 4.85) {\textbf{\color{orange!70!black}\faExpand\ margin}};

  % Border
  \fill[gray!40, rounded corners=3pt] (1,0.7) rectangle (7, 4.3);
  \draw[ultra thick, gray!70!black, rounded corners=3pt] (1,0.7) rectangle (7,4.3);
  \node[anchor=north west] at (1.1, 4.2) {\textbf{\color{gray!50!black}\faSquare\ border}};

  % Padding
  \fill[green!25, rounded corners=2pt] (1.3, 1) rectangle (6.7, 4);
  \node[anchor=north west] at (1.4, 3.9) {\textbf{\color{green!50!black}\faCompress\ padding}};

  % Content (innermost)
  \fill[blue!25, rounded corners=2pt] (2.2, 1.7) rectangle (5.8, 3.3);
  \node[anchor=north west] at (2.3, 3.2) {\textbf{\color{blue!60!black}\faFile*\ content}};
  \node at (4, 2.5) {\textit{textul, imaginea, ce e inauntru}};
\end{tikzpicture}
\end{center}
}

% =========================================================
%   CARD "ICON + TEXT" (utilitate pentru handout-uri)
% =========================================================
\newcommand{\iconcard}[3]{%
  \begin{tcolorbox}[
    enhanced,
    colback=partcolor!8,
    colframe=partcolor!30,
    boxrule=0.5pt,
    arc=6pt,
    left=12pt, right=12pt, top=8pt, bottom=8pt,
  ]
  \begin{minipage}[t]{0.08\linewidth}
    \color{partcolor}\fontsize{20}{20}\selectfont #1
  \end{minipage}%
  \begin{minipage}[t]{0.9\linewidth}
    \textbf{\color{partcolor!85!black}#2}\\
    #3
  \end{minipage}
  \end{tcolorbox}%
}
